\documentclass[11pt]{article}
\usepackage[T1]{fontenc}
\usepackage{lmodern,amsmath,amssymb,amsthm,booktabs}
\usepackage[margin=1in]{geometry}
\usepackage[colorlinks=true,linkcolor=blue,urlcolor=blue,citecolor=blue]{hyperref}
\usepackage{microtype,xurl}
\newtheorem{theorem}{Theorem}
\newtheorem{lemma}[theorem]{Lemma}
\newtheorem{corollary}[theorem]{Corollary}
\newcommand{\Z}{\mathbb Z}
\newcommand{\Q}{\mathbb Q}
\newcommand{\OK}{\mathcal O_K}
\newcommand{\blog}{https://terrytao.wordpress.com/2026/10/20/an-experimental-ai-math-challenge-using-tenth-discriminant-bounds-to-explore-arguments-ruling-out-siegel-zeroes/}
\title{An elementary determinant bound\\for the tenth discriminant\\
\large A proof with cutoff \(q<2^{15}\)}
\author{Research note prepared with Codex}
\date{Draft of 21 October 2026}
\begin{document}
\maketitle

\begin{abstract}
We prove that an imaginary quadratic field of class number one has
absolute fundamental discriminant \(q<32768\). A count of quadratic
norms supplies an auxiliary inert prime, at which a derivative
interpolation determinant is nonzero. Frobenius congruences at the
inert primes then force more divisibility than its complex absolute
value permits. The prime estimates are proved using factorials and
binomial coefficients, and the numerical argument consists of one
endpoint comparison. A separate integer sieve, using primes at most
31, excludes all remaining possibilities outside the nine classical
discriminants. The proof is independent of the earlier versions in
this project; its verification requires no numerical parameter search.
\end{abstract}

\section{Introduction and arithmetic preliminaries}

Tao's tenth-discriminant challenge \cite{tao} asks for explicit bounds
extracted from recent arguments concerning Landau--Siegel zeros,
including bounds that favor simple, human-verifiable proofs over
numerical strength. We give such a bound here. The smaller cutoff
\(q<570\) in \cite{v6} uses substantially more refinement; no result
from that note is an input below.

For a negative fundamental discriminant \(D\), put \(q=|D|\) and
\(K=\Q(\sqrt D)\).
\begin{theorem}\label{main}
If \(h(K)=1\), then \(q<32768=2^{15}\).
\end{theorem}

The argument uses standard quadratic ideal theory, finite fields,
Hadamard's inequality, and elementary analysis. All estimates specific
to the proof are established below. The determinant is never computed:
the supplementary code checks rational constants and a finite residue
sieve.

We begin with the classical inertness obstruction emphasized by
Silberman \cite{silberman}. Every nonrational element of \(\OK\) has
norm at least \(q/4\), by its imaginary part in the usual integral
basis. If \(h(K)=1\), a noninert rational prime \(p\) has a principal
integral ideal of norm \(p\). Its generator cannot be a rational
integer, whose norm would be a square. Consequently
\begin{equation}\label{small-inert}
 p<q/4\quad\Longrightarrow\quad p\text{ is inert}.
\end{equation}

For the proof of Theorem~\ref{main}, suppose henceforth that
\(h(K)=1\) and \(q\ge q_0:=32768\). Then \(q\) is odd, since
otherwise 2 ramifies. It is also prime: an odd fundamental
discriminant is squarefree, and a composite \(q\) would have a
ramified prime factor at most \(\sqrt q<q/4\). Inertness of 2 gives
\(q\equiv3\pmod8\). Thus the only ramified prime is \(q\), the units
are \(\pm1\). Put \(u=\sqrt q\). Then
\begin{equation}\label{norm-form}
 \OK=\Z+\Z\omega,\qquad \omega=(1+iu)/2,\qquad
 Q(a,b)=|a+b\omega|^2=(a+b/2)^2+qb^2/4.
\end{equation}
Write \(L=\log q\); all logarithms are natural.
We will use the elementary endpoint bounds
\begin{equation}\label{endpoint}
 \frac{207}{20}<\log q_0<\frac{21}{2},\qquad
 \sqrt{q_0}>180.
\end{equation}
Indeed, \(q_0=2^{15}\), \(69/100<\log2<7/10\), and
\(180^2<32768\). The function \((\log q)/\sqrt q\) decreases for
\(\log q>2\), so for every \(q\ge q_0\),
\begin{equation}\label{uniform-endpoint}
 L>207/20,\qquad \frac Lu<\frac7{120}.
\end{equation}

\section{Elementary prime estimates}\label{primes}

Write
\[
 \vartheta(x)=\sum_{p\le x}\log p,\qquad
 A(x)=\sum_{p\le x}\frac{\log p}{p},\qquad
 B=\sum_p\frac{\log p}{p^2}.
\]
A letter \(p\) under a sum always denotes a rational prime. The
following Chebyshev-style estimates suffice; no lower bound for
\(\vartheta\) is required.

\begin{lemma}\label{primelemma}
For \(x\ge1\) and integers \(m\ge2\),
\begin{equation}\label{prime-estimates}
 \vartheta(x)<\frac32x,\qquad
 \log m-\frac73<A(m)<\log m+\frac32,\qquad B<\frac23.
\end{equation}
\end{lemma}
\begin{proof}
Write \(n=2k\) or \(n=2k+1\). Every prime in
\((\lceil n/2\rceil,n]\) divides the central binomial coefficient
\(\binom n{\lfloor n/2\rfloor}\). This coefficient is at most
\(2^{2k}\): in the odd case there are two equal central coefficients.
Thus
\[
 \vartheta(n)-\vartheta(\lceil n/2\rceil)
 \le 2\lfloor n/2\rfloor\log2.
\]
Starting at \(n=1\), induction gives \(\vartheta(n)<2n\log2\).
Taking \(n=\lfloor x\rfloor\), and using \(\log2<7/10<3/4\),
proves the first assertion.

Every prime other than 2 is odd, and \(f(x)=\log x/x^2\) decreases
for \(x\ge3\). For each odd integer \(n\ge5\),
\(f(n)\le\tfrac12\int_{n-2}^n f(x)\,dx\). Hence
\begin{align*}
 B&\le\frac{\log2}{4}+\frac{\log3}{9}
                  +\frac12\int_3^\infty\frac{\log x}{x^2}\,dx\\
  &=\frac{\log2}{4}+\frac5{18}\log3+\frac16
    <\frac{233}{360}<\frac23,
\end{align*}
using \(\log3<11/10\).

The factorial valuation formula and integral comparison give
\[
 m\log m-m<\log(m!)
 \le mA(m)+m\sum_p\frac{\log p}{p(p-1)}
 \le mA(m)+2mB<mA(m)+\frac43m.
\]
This proves the lower bound for \(A(m)\). In the other direction,
\(\log(m!)\ge mA(m)-\vartheta(m)\) and
\(\log(m!)\le m\log m\) prove the upper bound.
\end{proof}

\section{Norm counts and an auxiliary inert prime}\label{norms}

Pratt \cite{pratt} suggested using class number one to bound split
primes by norm representations, then using the resulting supply of
inert primes in a determinant argument. We implement this with an
unsieved norm count and a weighted interval of ordinary primes.

Let \(\pi_+(x)\) count split rational primes at most \(x\). For each
integer \(b>0\), the real interval of \(a\) satisfying
\(Q(a,b)\le x\) has length \(2\sqrt{x-qb^2/4}\), when nonempty.
Bounding the number of integers by length plus 1, and the decreasing
length sum by its integral, gives
\[
 \#\{(a,b)\in\Z\times\Z_{>0}:Q(a,b)\le x\}
 \le\frac{\pi x}{u}+\frac{2\sqrt x}{u}.
\]
Each split prime has two generators with \(b>0\), one for each prime
ideal above it: each ideal has generators \(\pm\alpha\), exactly one
of which has positive \(b\). Consequently
\begin{equation}\label{split-count}
 \pi_+(x)\le\frac{\pi x}{2u}+\frac{\sqrt x}{u}.
\end{equation}

For \(X\ge Y\ge q/4\), partial summation, followed by
\eqref{split-count}, yields
\begin{align}
 \sum_{\substack{Y<p\le X\\p\text{ split}}}\frac{\log p}{p}
 &\le\frac{\pi}{4u}
       \bigl((\log X)^2-(\log Y)^2+2\log Y\bigr)
       +\frac{2(\log Y+1)}{u\sqrt Y}.\label{interval-weight}
\end{align}
To see the boundary terms explicitly, start with
\[
 \frac{\log X}{X}\pi_+(X)-\frac{\log Y}{Y}\pi_+(Y)
 +\int_Y^X\pi_+(t)\frac{\log t-1}{t^2}\,dt.
\]
Discard the nonpositive lower boundary term. Substitution of
\eqref{split-count} then gives the right side of
\eqref{interval-weight}, together with the further nonpositive term
\(-(\log X+2)/(u\sqrt X)\).

Put \(S_+(X)=\sum_{p\le X,\ p\text{ split}}\log p/p\).
There are no split primes at or below \(q/4\), by
\eqref{small-inert} and the fact that \(q/4\) is not an integer.
Taking \(Y=q/4\) in \eqref{interval-weight}, and using
\(\log(q/4)>2\) and \(\pi<16/5\), gives
\begin{equation}\label{split-weight}
 S_+(X)\le\frac4{5u}(\log X)^2+\frac{4(L+1)}q
 \qquad(X\ge q/4).
\end{equation}

\begin{lemma}\label{auxiliary}
There is an inert prime \(R\) with \(q^4<R\le q^5\).
\end{lemma}
\begin{proof}
Lemma~\ref{primelemma} shows that the total weight \(\log p/p\)
of primes in \((q^4,q^5]\) is greater than \(L-23/6\).
By \eqref{interval-weight}, the split-prime weight there is at most
\[
 \frac4{5u}(9L^2+8L)+\frac{2(4L+1)}{u q^2}.
\]
It is therefore enough to check
\begin{equation}\label{aux-inequality}
 \frac{23}{6L}+\frac4{5u}(9L+8)
             +\frac{2(4+1/L)}{u q^2}<1.
\end{equation}
Each term decreases as \(q\) increases in our range. For the middle
term, differentiation with respect to \(L\) gives a factor
\(9-(9L+8)/2<0\); the other terms plainly decrease.
At \(q=q_0\), \eqref{endpoint} bounds the left side by
\[
 \frac{10}{27}+\frac{41}{90}
       +\frac{2(4+20/207)}{180\cdot32768^2}
 <\frac{10}{27}+\frac{41}{90}+\frac1{1000}<\frac56.
\]
The only ramified prime, \(q\), is below the interval. Since the
split primes cannot account for its full weight, an inert prime exists.
\end{proof}

\section{The derivative determinant}\label{determinant}

We use Green's derivative matrix \cite[\S2]{green}, with a
rectangular set of algebraic integers. Choose \(R\) as in
Lemma~\ref{auxiliary} and put
\[
 S=\lfloor R/12\rfloor,\qquad J=\lceil R/48\rceil,\qquad
 N=RS,\qquad T=2\sqrt{Ru}.
\]
Take all points
\[
 z=a+b\omega,\qquad
 |a|\le A_0:=\lceil\sqrt{Ru}\rceil,\qquad
 |b|\le B_0:=\lceil\sqrt{R/u}\rceil.
\]
Their number is \(n=(2A_0+1)(2B_0+1)\ge4R\), so \(nJ\ge RS\).
Since \(|\omega|\le3u/5\), \(u>1\), and \(R>q^4>100u\),
\[
 |z|\le\frac85\sqrt{Ru}+1+\frac35u
       \le\frac{44}{25}\sqrt{Ru}<T.
\]
Here \(\sqrt{Ru}>10u\). Also \(2T<R\) and \(1<T<R\).

Form the matrix whose rows are indexed by \((r,s)\),
\(0\le r<R\), \(0\le s<S\), and whose columns are indexed by a
point \(z_i\) and an order \(0\le k<J\). The entry is
\begin{equation}\label{matrix-entry}
 \left.\partial_z^k(z^r w^s)\right|_{(z_i,\bar z_i)}.
\end{equation}
These are ordinary derivatives, without division by factorials.

\subsection{Nonvanishing over a finite field}

The following argument is the auxiliary-prime method of van Dijk
\cite{vandijk}, highlighted in Green's comment \cite{greencomment}.
Green also reports essentially the same method in unpublished work
of Jacob Pescador; the present note uses the public argument and does
not rely on that manuscript.

Reduce \eqref{matrix-entry} modulo the inert prime ideal \(R\OK\).
The residue field has \(R^2\) elements. The nodes remain distinct:
every nonzero element of \(R\OK\) has modulus at least \(R\),
because a nonzero algebraic integer has positive integral norm,
whereas any difference of two nodes has modulus less than \(2T<R\).
Frobenius gives \(\bar z_i=z_i^R\) in this field.

If a nonzero row combination vanished, let \(F(z,w)\) be its
polynomial and substitute \(w=z^R\). The exponents \(r+Rs\) are
distinct, so \(F(z,z^R)\) is nonzero and has degree at most \(RS-1\).
Its derivatives agree with the prescribed partial derivatives,
because the derivative of \(z^R\) is zero in characteristic \(R\).
Since \(J\le R\) and \(k!\ne0\) for \(0\le k<J\), each node is
a root of multiplicity at least \(J\). This contradicts \(nJ\ge RS\).

The matrix therefore has full row rank. Fix one nonzero \(N\)-square
minor \(D_0\in\OK\), and set
\[
 \Delta=D_0\overline{D_0}\in\Z_{>0}.
\]
Only this minor is used from now on.

\subsection{The complex absolute-value bound}

Each entry in row \((r,s)\) has modulus at most
\(T^{r+s}(R/T)^J\): the derivative contributes at most \(R^k\),
and an entry with \(k>r\) is zero. Hadamard's inequality gives
\[
 \frac{\log\Delta}{R^2S}
 \le\frac{\log N}{R}+\left(1+\frac SR\right)\log T
                           +\frac{2J}{R}\log(R/T).
\]
Now \(S/R\le1/12\), \(J/R\le1/48+1/R\), and \(N<R^2\).
Thus the coefficient of \(\log T\), after combining the main
terms, is \(25/24\). Since
\(\log T=\log2+\tfrac12\log R+\tfrac14L\), we obtain
\begin{equation}\label{upper-det}
 \frac{\log\Delta}{R^2S}
 \le\frac9{16}\log R+\frac{25}{96}L
       +\frac{25}{24}\log2+\frac{4\log R}{R}.
\end{equation}

\subsection{Divisibility at the inert primes}

We use Green's triangular Frobenius transformation
\cite[Lemma~5.1]{green}, the stronger divisibility input suggested
by Pratt \cite{pratt}. The argument below applies also to \(p=2\).
For an inert prime \(p\), put
\[
 h_{r,s}=\min(\lfloor r/p\rfloor,S-1-s)
\]
and replace each row polynomial \(z^rw^s\) by
\[
 z^{r-ph_{r,s}}(z^p-w)^{h_{r,s}}w^s.
\]
The leading term is \(z^rw^s\), and every other term has smaller
\(z\)-degree while remaining in the rectangle of row indices.
Ordering rows by \(z\)-degree makes this an integer triangular
transformation of determinant one.

At each node \(z_i^p-\bar z_i\in p\OK\), by Frobenius. Every
positive ordinary derivative of \(z^p-w\) is also divisible by
\(p\). By the product rule, all entries in the transformed row
are therefore divisible by \(p^{h_{r,s}}\). Taking norms gives
\begin{equation}\label{valuation}
 v_p(\Delta)\ge2\sum_{r,s}h_{r,s}.
\end{equation}
The row transformation may depend on \(p\); each divisibility is
nevertheless a statement about the same integer \(\Delta\).

We bound this sum with generous constants. Direct summation gives,
for every integer \(k\ge0\),
\[
 \sum_{s<S}\min(k,S-1-s)\ge Sk-\frac{k(k+1)}2.
\]
Writing \(R=ap+b\), \(0\le b<p\), gives
\[
 \sum_{r<R}\lfloor r/p\rfloor
 =\frac{R(R-p)+b(p-b)}{2p}\ge\frac{R(R-p)}{2p}.
\]
Also \(\lfloor r/p\rfloor(\lfloor r/p\rfloor+1)
\le r^2/p^2+r/p\). Substituting these estimates yields
\begin{equation}\label{valuation-bound}
 \frac{2\sum_{r,s}h_{r,s}}{R^2S}
 \ge\left(1-\frac{R-1}{2RS}\right)\frac1p-\frac1R
                  -\frac{(R-1)(2R-1)}{6RS p^2}.
\end{equation}
Let \(I=\sum_{p\le R,\ p\text{ inert}}\log p/p\). Summing
\eqref{valuation-bound}, weighted by \(\log p\), and using
\(S\ge R/13\) and Lemma~\ref{primelemma}, we find
\begin{align}
 \frac{\log\Delta}{R^2S}
 &\ge I-\frac{I}{2S}-\frac{\vartheta(R)}R-\frac{R}{3S}B\notag\\
 &>I-\frac{I}{2S}-\frac32-\frac{26}{9}
  =I-\frac{I}{2S}-\frac{79}{18}.\label{lower-det}
\end{align}
Here \(S\ge R/12-1\ge R/13\), since \(R>156\).

\section{One endpoint comparison}

Separating the three splitting types in \(A(R)\), and noting that
the only ramified contribution is \(L/q\), gives
\[
 I>\log R-\frac73-S_+(R)-\frac Lq,
 \qquad I<\log R+\frac32.
\]
Combine this with \eqref{upper-det}, \eqref{lower-det}, and
\eqref{split-weight}. If
\[
 G(x)=\frac7{16}x-\frac{25}{96}L-\frac4{5u}x^2,
\]
the resulting necessary inequality is
\begin{equation}\label{G-upper}
 G(\log R)<\frac{121}{18}+\frac{25}{24}\log2+E,
 \qquad
 E=\frac{5L+4}q+\frac{I}{2S}+\frac{4\log R}{R}.
\end{equation}
All boundary and rounding terms are included in \(E\). In fact
\begin{equation}\label{error-bound}
 E<\frac{5L+4}q+\frac{(21/2)\log R+39/4}{R}
 <\frac{113}{65536}+\frac{1719}{4\cdot10^{16}}<\frac1{500}.
\end{equation}
For the second inequality, the first function decreases with \(q\)
in our range, so use \eqref{endpoint} at \(q=q_0\).
The second decreases with \(R\); use \(R>q_0^4>10^{16}\) and
\(\log(10^{16})<40\). Thus \(\log2<7/10\) gives
\begin{equation}\label{G-simple}
 G(\log R)<\frac{121}{18}+\frac{35}{48}+\frac1{500}
 <\frac{15}{2}.
\end{equation}

On the other hand, \(4L<\log R\le5L\) by
Lemma~\ref{auxiliary}, and \(G\) is concave. Its minimum on
\([4L,5L]\) is therefore at an endpoint. By
\eqref{uniform-endpoint},
\[
 G(5L)-G(4L)
 =L\left(\frac7{16}-\frac{36}{5}\frac Lu\right)
 >\frac7{400}L>0.
\]
It follows that
\[
 G(\log R)\ge G(4L)
 =\frac{143}{96}L-\frac{64}{5}\frac{L^2}{u}.
\]
Together with \eqref{G-simple}, this forces
\[
 \frac{143}{96}
 <\frac{64}{5}\frac Lu+\frac{15}{2L}
 <\frac{56}{75}+\frac{50}{69}
 =\frac{846}{575}<\frac{143}{96},
\]
a contradiction. This proves Theorem~\ref{main}.

\section{The remaining finite range}\label{finite}

The residue sieve is the elementary small-prime method discussed by
Pratt \cite{prattsieve}. For \(q>8\), \eqref{small-inert} at 2
forces \(q\equiv3\pmod8\), without requiring the large-discriminant
assumption used above. For each odd prime \(p<q/4\), inertness
requires \(-q\) to be a nonzero nonsquare modulo \(p\).

There are 4095 integers \(q\equiv3\pmod8\) with
\(8<q<32768\). Test all of them, including nonfundamental values,
against the primes
\[
 3,5,7,11,13,17,19,23,29,31.
\]
The following complete integer computation returns exactly the list
shown in its assertion:
\begin{verbatim}
primes = (3, 5, 7, 11, 13, 17, 19, 23, 29, 31)
squares = {p: {a*a % p for a in range(p)} for p in primes}
survivors = [q for q in range(11, 32768, 8)
             if all(4*p >= q or (-q) % p not in squares[p]
                    for p in primes)]
assert survivors == [11, 19, 43, 67, 163]
\end{verbatim}
The negative fundamental discriminants with \(q\le8\) have
absolute values \(3,4,7,8\). We have therefore proved, with the
finite computation just specified:
\begin{corollary}
If \(D<0\) is a fundamental discriminant and \(h(D)=1\), then
\[
 |D|\in\{3,4,7,8,11,19,43,67,163\}.
\]
\end{corollary}
As an additional check of the familiar converse, the supplied script
enumerates primitive reduced positive definite forms of each of these
discriminants. Using \(|b|\le a\le c\) and \(b\ge0\) when
\(|b|=a\) or \(a=c\), it finds one form in each case. This uses the
classical correspondence between reduced forms and ideal classes;
it is separate from the exclusion argument.

\section*{Reproducibility}

The supplementary files use only the Python standard library. Run
\begin{verbatim}
python3 verify_constants.py
python3 verify_finite.py
\end{verbatim}
The first checks the numerical inequalities by rational arithmetic.
For logarithms it uses 16 terms of
\[
 \log y=2\sum_{j\ge0}\frac{t^{2j+1}}{2j+1},\qquad
 t=\frac{y-1}{y+1},\qquad 1\le y\le2,
\]
with tail at most \(2t^{33}/[33(1-t^2)]\), after reduction by powers
of 2. The elementary identity
\(\pi/4=\arctan(1/2)+\arctan(1/3)\), with three terms of each
alternating series, verifies \(\pi<16/5\). These also verify
\(69/100<\log2<7/10\), \(\log3<11/10\), and \(\log10<5/2\),
the only logarithmic inputs to the proof.

The second script performs the exhaustive residue check, verifies
its residue tables by Euler's criterion, and enumerates the nine
reduced-form lists. The scripts and their recorded outputs accompany
this note. The monotonicity arguments in the text cover the infinite
range; the scripts check its endpoints, not samples of that range.
These computations are not formal verification, and independent
mathematical review of this draft is pending.

\section*{Acknowledgements and AI disclosure}

The author thanks the participants in Tao's discussion for making
their arguments and suggestions available. Their specific mathematical
contributions are cited where used above. This note, including its
derivation, exposition, and verification code, was developed with
OpenAI Codex. Green \cite{green} credits ChatGPT for the mathematical
arguments in his source paper; van Dijk \cite{vandijk} credits
GPT-6 Astra and Opus 5.5 for his new arguments.

\clearpage
\begin{thebibliography}{9}

\bibitem{green}
B.~Green, \emph{On Landau--Siegel zeros}, version 2 (2026),
\href{https://hexagonmath.org/2610.00311v2}{hexagon:2610.00311v2}.

\bibitem{greencomment}
B.~Green, Comment on Tao's tenth-discriminant challenge,
20 October 2026, 12:07 pm,
\href{\blog\#comment-698463}{comment 698463}.

\bibitem{prattsieve}
K.~Pratt, Comment on small-prime congruences and finite sieving,
20 October 2026, 11:15 am,
\href{\blog\#comment-698460}{comment 698460}.

\bibitem{pratt}
K.~Pratt, Comment on norm counts, auxiliary inert primes, and
determinant divisibility, 20 October 2026, 1:58 pm,
\href{\blog\#comment-698476}{comment 698476}.

\bibitem{silberman}
L.~Silberman, Comment on small-prime inertness and weighted prime
sums, 20 October 2026, 8:49 am,
\href{\blog\#comment-698449}{comment 698449}.

\bibitem{tao}
T.~Tao, \emph{An experimental AI-math challenge: using tenth
discriminant bounds to explore arguments ruling out Siegel zeroes},
20 October 2026, \href{\blog}{What's new}.

\bibitem{v6}
\emph{A smoothed Frobenius determinant bound for the
tenth discriminant}, version 0.6.0 (2026),
\href{https://github.com/twinheterodox/tenth-discriminant/releases/tag/v0.6.0}{GitHub release}.

\bibitem{vandijk}
R.~van Dijk, \emph{A Frobenius determinant proof of the non-existence
of Landau--Siegel zeros}, version 1 (2026),
\href{https://hexagonmath.org/2610.00409v1}{hexagon:2610.00409v1}.

\end{thebibliography}
\end{document}
